%% ma: L12-B15-0010
%% lop: 12
%% chuong: 5
%% bai: 15
%% dang: 12.15.8
%% loai: DS
%% muc: [NB, TH, VD, VDC]
%% dap_an: "SĐSĐ"
%% nguon: "(NBV)-ĐỀ SỐ 14-MỖI NGÀY 1 ĐỀ THI 2025.pdf (D14, MỖI NGÀY 1 ĐỀ - Đề số 14), Phần II Câu 4"
%% kien_thuc: "Bài 8 (tích vô hướng, độ dài vectơ), Bài 15 (phương trình đường thẳng); quỹ tích điểm cách một đường thẳng cho trước một khoảng không đổi (mặt trụ)"
%% trang_thai: cho-duyet
%% ly_do: "Dạng mới 12.15.8 đề xuất, chờ giáo viên duyệt"
%% ghi_chu: "Đáp án gốc đều đúng. Đề không có hình."
\begin{ex}
Trong không gian $Oxyz$, cho các điểm $A(1;1;0)$, $B(5;-3;2)$ và $C(0;4;-1)$. Xét các điểm $M$ thay đổi trong không gian sao cho diện tích tam giác $ABM$ bằng $6\sqrt2$.
\choiceTF
{Đoạn thẳng $AB$ có độ dài bằng $3$.}
{\True Đường thẳng $AB$ có phương trình là $\dfrac{x-1}{2}=\dfrac{y-1}{-2}=\dfrac z1$.}
{Khoảng cách từ điểm $C$ tới đường thẳng $AB$ bằng $2\sqrt2$.}
{\True Đoạn thẳng $MC$ có độ dài nhỏ nhất bằng $\sqrt2$.}
\loigiai{$\overrightarrow{AB}=(4;-4;2)$, $AB=\sqrt{16+16+4}=6$.

a) $AB=6\ne3$. Sai.

b) $\overrightarrow{AB}=2(2;-2;1)$ nên $\vec u=(2;-2;1)$ là vectơ chỉ phương, đường thẳng $AB$ qua $A(1;1;0)$ có phương trình $\dfrac{x-1}{2}=\dfrac{y-1}{-2}=\dfrac z1$. Đúng.

c) $\overrightarrow{AC}=(-1;3;-1)$, $AC^2=11$; hình chiếu của $\overrightarrow{AC}$ lên $\vec u$ có độ dài $\dfrac{|\overrightarrow{AC}\cdot\vec u|}{|\vec u|}=\dfrac{|-2-6-1|}{3}=3$. Vậy $d(C,AB)=\sqrt{11-9}=\sqrt2\ne2\sqrt2$. Sai.

d) $S_{ABM}=\dfrac12AB\cdot d(M,AB)=6\sqrt2\Rightarrow d(M,AB)=2\sqrt2$, nên $M$ nằm trên mặt trụ trục $AB$, bán kính $2\sqrt2$. Điểm $C$ có $d(C,AB)=\sqrt2<2\sqrt2$ nằm trong mặt trụ nên $MC_{\min}=2\sqrt2-\sqrt2=\sqrt2$. Đúng.}
\end{ex}
